\clearpage\begingroup\SourceLeaf{687b}
\ArithmeticHead{52}{L}{H. GEPPERT, ÜBER GAUSS’ ARBEITEN ZUR MECHANIK UND POTENTIALTHEORIE.}
\makeatletter\tagsleft@true\let\veqno\@@leqno\makeatother
\renewcommand{\thefootnote}{\arabic{footnote})}\setcounter{footnote}{102}
\begin{center}\textbf{V. Arbeiten zur praktischen Mechanik.}\end{center}
\noindent\hspace*{\parindent}Die praktischen Arbeiten über den Erdmagnetismus und die Konstruktion
der benötigten Apparate haben bei \textsc{Gauss} auch einige in das Gebiet der
praktischen Mechanik gehörige Untersuchungen ausgelöst; diese sind vor allem
drei: 1. die genaue Bestimmung der Schwingungsdauer einer Magnetnadel,
2. die Ermittlung der Torsion eines Fadens, 3. die Bestimmung des Trägheitsmomentes
einer Magnetnadel.

\begin{center}\textls[100]{1. Schwingungsdauer einer Magnetnadel.}\end{center}
\noindent\hspace*{\parindent}In der \textit{Intensitas vis magneticae}, die \textsc{Gauss} am 15. Dez. 1832 der Göttinger
Sozietät vorlegte, hatte er sich die Aufgabe gestellt, die Intensität der
erdmagnetischen Kraft zu ermitteln. Er führte dieselbe auf die Bestimmung
der Schwingungsdauer einer gegebenen Magnetnadel zurück\footnote{Vergl. hierzu den Essay von \textsc{Cl. Schaefer}, \textit{Über Gauss’ physikalische Arbeiten}, Werke XI, 2.}, und im wesentlichen
beruhen alle Methoden der \textit{Intensitas} auf der genauen Bestimmung
dieser Schwingungsdauer, auf die also grosse Präzision zu verwenden ist. Die
Auffindung des magnetischen Meridians erfordert überdies die Bestimmung
der Ruhelage der schwingenden Nadel. Mit diesen beiden Aufgaben befassen
sich drei besondere Abhandlungen von \textsc{Gauss}: a) \textit{Das in den Beobachtungsterminen
anzuwendende Verfahren}, 1836\footnote{Resultate aus den Beobb. des magnetischen Vereins, 1836, S. 34—50. Werke V, S. 541—556.}, b) \textit{Anleitung zur Bestimmung der
Schwingungsdauer einer Magnetnadel}, 1837\footnote{Resultate aus den Beobb. des magnetischen Vereins, 1837, S. 58—80. Werke V, S. 374—394.}, c) \textit{Über ein Mittel, die Beobachtung
von Ablenkungen zu erleichtern}, 1839\footnote{Resultate aus den Beobb. des magnetischen Vereins, 1839, S. 52—62. Werke V, S. 395—403.}.
\clearpage\SourceLeaf{688}\ArithmeticHead{53}{R}{V. ARBEITEN ZUR PRAKTISCHEN MECHANIK.}

Die Bestimmung der Ruhelage einer schwingenden Nadel kann erfolgen
entweder aus den abgelesenen aufeinanderfolgenden maximalen Elongationen
\(a,b,c,\, .\, .\, .\) oder, nach Kenntnis der Schwingungsdauer, aus der Ablesung zweier
um eine Schwingungsdauer unterschiedener Stellungen. Beide Verfahren entwickelt
\textsc{Gauss} in der Abhandlung a). Will man der durch die Dämpfung
bedingten Abnahme des Schwingungsbogens Rechnung tragen, und ist \(\theta=e^{-\lambda}\)
der Dämpfungsfaktor, also \(\lambda\) das logarithmische Dekrement (dieser Name selbst
rührt erst von \textsc{Gauss} her), so wird die Ruhelage statt durch \(\frac{a+b}{2}\) durch
\(\frac14(a+2b+c)\) ausgedrückt; der dann begangene Fehler ist, wie \textsc{Gauss} im Briefe
vom 20. Okt. 1837 an \textsc{Schumacher}\footnote{\textsc{Gauss-Schumacher}-Briefwechsel, 3, S. 180—182.} errechnet, gleich \(\frac{(a-c)^2}{4(a+2b+c)}\), also zu
vernachlässigen. Will man genauer, was bei starken Deklinationsänderungen
erforderlich ist, die Ruhelage in einem gegebenen Augenblicke \(t_0\) bestimmen,
so muss man vorher die Schwingungsdauer \(T\) der Nadel kennen und die Elongationen
zu den Zeiten \(t_0-\frac12T\) und \(t_0+\frac12T\) messen; deren Mittel ergibt die
gesuchte Ruhelage, solange es sich um eine reine Sinusschwingung handelt.
Dieses Verfahren erlaubt insbesondere, die unexakte Messung an den Stellen
grösster Elongation durch die viel genauere an der Stelle grösster Geschwindigkeit
zu ersetzen, ein Prinzip, das \textsc{Gauss} bereits 1829 besass\footnote{Vergl. \textsc{Schaefer}, a.\,a.\,O. S. 7—8.}. Mit der vollständigen
Untersuchung des Einflusses der Dämpfung auf Schwingungsdauer
und Ruhelage beschäftigt sich \textsc{Gauss} in der Abhandlung b). Die Aufgabe ist
gleichbedeutend mit der Integration der Differentialgleichung
\[
\ddot{x}+2\varepsilon\dot{x}+n^2(x-p)=0,\tag{1}
\]
\noindent wenn die Dämpfung proportional der Geschwindigkeit ist und \(p\) die Ruhelage
bedeutet. Das allgemeine Integral
\[
x=p+Ae^{-\varepsilon t}\sin\sqrt{n^2-\varepsilon^2}(t-B)
\]
\noindent zeigt, dass Amplituden und Schwingungsdauer durch die Dämpfung affiziert
werden. Setzt man \(\frac\varepsilon n=\sin\varphi\), so gilt zwischen der gedämpften und der ungedämpften
Schwingungsdauer \(T'\) bezw. \(T\) die Relation \(T=T'\cos\varphi\), und der
Dämpfungsfaktor der Amplituden wird:
\[
\theta=e^{-\varepsilon T'}=e^{-\pi\operatorname{tang}\varphi},
\]
\clearpage\SourceLeaf{689}\ArithmeticHead{54}{L}{H. GEPPERT, ÜBER GAUSS’ ARBEITEN ZUR MECHANIK UND POTENTIALTHEORIE.}
\noindent also das logarithmische Dekrement: \(\lambda=\pi\operatorname{tang}\varphi\). Die Kenntnis von \(\lambda\) vermittelt
somit die Reduktion des gemessenen \(T'\) auf \(T\).

Um \(T'\) zu messen, beobachtet man die Durchgangszeiten an einem fixierten
Skalenpunkte; bei ungedämpfter Bewegung ergäbe deren Mittel die Zeit der
dazwischenliegenden maximalen Elongation; bei gedämpfter Bewegung hingegen
ist hieran eine Korrektion \(\delta\) anzubringen, die sich aus einer transzendenten
Gleichung berechnet und bei Beobachtung in der Umgebung des Ruhepunktes
gleich \(-\frac{T'}{\pi}\varphi\) ist. Diese Korrektion verschiebt also alle Elongationszeiten
um ein konstantes Stück, spielt somit für deren Unterschied, d.\,h. \(T'\), keine
Rolle. Auch die Berechnung des Ruhepunktes \(p\) modifiziert sich infolge der
Dämpfung: Sind \(x\) und \(x'\) die Nadelstellungen zu den Zeiten \(t\) und \(t+T'\), so ist
\[
p=x+\frac{x'-x}{1+\theta}.\tag{2}
\]
\noindent Schliesslich gelten die entwickelten Formeln nur solange, als es sich um infinitesimale
Schwingungsbogen handelt. Ist hingegen der Schwingungsbogen
\(G\) endlich und die zugehörige Schwingungsdauer \(\widetilde T\), so findet sich die zu unendlichkleinen
Schwingungsbogen gehörende Schwingungsdauer aus der Näherungsformel
\[
T=\widetilde T\left(1-\frac{G^2}{64}\right),
\]
\noindent die durch Entwicklung eines elliptischen Integrals entsteht. Auch diese Korrektionsformel
ändert sich ab, wenn man es mit gedämpften Schwingungen zu
tun hat.

\textsc{Gauss} geht noch einen Schritt weiter und untersucht die Schwingungen,
die entstehen, wenn die sie hervorrufende Kraft, also im vorliegenden Falle
die Deklination, sich während der Versuche beträchtlich ändert. Dies kommt
auf die Lösung der Differentialgleichung
\[
\ddot{x}+2\varepsilon\dot{x}+n^2(x-p-at)=0\tag{3}
\]
\noindent hinaus, deren Integral
\[
x=p-\frac{2a\varepsilon}{n^2}+at+Ae^{-\varepsilon t}\sin\sqrt{n^2-\varepsilon^2}(t-B)
\]
\noindent zeigt, dass nur eine Verschiebung des Ruhepunktes eintritt. Diese kann man
in der Weise kenntlich machen, dass der durch (2) bestimmte Ruhepunkt zu
der Zeit \(t+\frac{T'}{1+\theta}-\frac{T'}{\pi}\sin2\varphi\) gehört. Auf diese Korrektion bezieht sich das
\clearpage\SourceLeaf{690}\ArithmeticHead{55}{R}{V. ARBEITEN ZUR PRAKTISCHEN MECHANIK.}
\noindent Werke XI, 1, S. 67—68 abgedruckte Nachlassstück aus dem Handbuch 19,
Be, das auf die Zeit zwischen 1834 und 1837 zu datieren ist\footnote{Vergl. \textsc{Schaefer} a.\,a.\,O. S. 131.}.

Die \textsc{Gauss}sche Abhandlung b) setzt sich also das Ziel, die theoretischen Erkenntnisse
über die Differentialgleichungen (1) und (3) in die Versuchspraxis
umzusetzen; sie ist zugleich die erste Arbeit, die dies leistet. \textsc{Poisson}\footnote{\textsc{S. D. Poisson}, \textit{Traité de mécanique}, 2. éd. 1833, vol. 1, p. 349 ff.} hatte
sich 1833 mit der Gleichung (1) beschäftigt, den Charakter der Lösung untersucht
und, ebenso wie \textsc{Gauss}, auf die zweifache Möglichkeit einer aperiodischen
und quasiperiodischen Lösung hingewiesen, es aber dabei bewenden
lassen. Zu dem auch heutigentags allgemein üblichen \textsc{Gauss}schen Verfahren
sei ergänzend bemerkt, dass \textsc{Gauss} eine Korrektion nicht berücksichtigt hat,
auf die \textsc{Lamont}\footnote{\textsc{J. Lamont}, \textit{Reduction der Schwingungen eines Magneten auf den luftleeren Raum.} \textsc{Poggendorfs} Ann. d. Phys. 71, 1847, S. 124—128.} aufmerksam machte, d.\,i. die Korrektion auf den luftleeren
Raum, die nötig wird, weil die schwingende Nadel auch die umgebende Luft
in Mitschwingung versetzt und infolgedessen ihre Schwingungsdauer vermindert.
In gewöhnlichen Fällen beträgt diese Korrektion nach \textsc{Lamont} etwa \(\frac1{8330}T\),
um diesen Betrag ist \(T\) zu vergrössern.

Die dritte \textsc{Gauss}sche Abhandlung c) gibt die zweckmässigste Bestimmung
einer die Magnetnadel ablenkenden Kraft. Modifiziert letztere weder die
Schwingungsdauer noch die Dämpfung, so verändert sie nur die Ruhelage der
schwingenden Nadel. \textsc{Gauss} bestimmt letztere folgendermassen: Die Versuchszeit
wird in vier Abschnitte eingeteilt, die Epoche a) vor dem Zeitpunkt \(t'\),
b) zwischen \(t'\) und \(t''\), c) zwischen \(t''\) und \(t'''\), d) nach \(t'''\). Während der Zeiträume
a) und c) wirkt die zusätzliche Kraft nicht, dagegen lässt man sie
während b) und d) voll einwirken. Dann ist die Aufgabe zu lösen, durch zweckmässige
Wahl der Zeitintervalle zu erreichen, dass bis auf den veränderten
Ruhepunkt die Schwingung in der Epoche d) genau die gleiche ist wie die
in a). Ist \(T\) die ungedämpfte Schwingungsdauer, und setzt man \(t''-t'=qT\),
\(t'''-t''=rT\), so erfordert diese Aufgabe die Lösung zweier transzendenter
Gleichungen für \(q\) und \(r\). Liegt keine Dämpfung vor, so findet sich \(q=r=\frac13\),
also ein einfaches Drittelungsverfahren; liegt dagegen Dämpfung vor, so wird
mit zunehmendem Dekrement \(q\) grösser und \(r\) kleiner. \textsc{Gauss} gibt die
\clearpage\SourceLeaf{691}\ArithmeticHead{56}{L}{H. GEPPERT, ÜBER GAUSS’ ARBEITEN ZUR MECHANIK UND POTENTIALTHEORIE.}
\noindent tabellarische Lösung der transzendenten Gleichung an. Auf diese Methode bezieht
sich auch ein Werke XI, 1, S. 62—63 abgedrucktes Nachlassstück aus dem
Handbuch 15, Ba, vom Jahre 1838\footnote{Vergl. \textsc{Schaefer}, a.\,a.\,O. S. 133.}.

Ergänzend wollen wir noch bemerken, dass \textsc{Gauss} 1835 sich theoretisch
und praktisch auch mit der Theorie der erzwungenen Schwingungen und dem
Phänomen der Resonanz beschäftigt hat, über die er in dem aufschlussreichen
Briefe an \textsc{Olbers} vom Februar 1835\footnote{\textsc{Gauss-Olbers}-Briefwechsel, 2, S. 617—620.} berichtet.
\clearpage\endgroup
